% =========================================================================
% Chapter 2: Review of Literature
% Chronological: International -> National (India) -> Regional (J&K)
% =========================================================================

\chapter{Review of Literature}\label{ch:litreview}

The literature is arranged by date, oldest first, in three geographic tiers. International work on \textit{Dipsacus}, related Caprifoliaceae, plant antimicrobials, and Gram-negative envelopes comes first. Indian work on medicinal-plant extraction and clinical Gram-negative resistance comes second. Kashmir Himalayan ethnobotany and laboratory studies of \taxa{Dipsacus inermis} come last.

\section{International studies}

\citet{jensen1979cantleyoside} isolated and characterised novel bis-iridoid glucosides from \taxa{Dipsacus sylvestris}, establishing cantleyoside as a dimeric iridoid glucoside formed through bridge condensation between iridoid and secoiridoid units. The investigation demonstrated that cyclopentanoid monoterpene iridoids, including loganin and sweroside, constitute distinctive chemotaxonomic markers across the order Dipsacales. Owing to their polyhydroxylated glucopyranosyl substituents, these iridoid glucosides possess high water solubility and partition readily into boiling aqueous decoctions and Soxhlet liquors.

\citet{li1995role} identified the tripartite organisation of the primary resistance-nodulation-division (RND) multidrug efflux system in \taxa{Pseudomonas aeruginosa}, encoded by the \textit{mexA-mexB-oprM} operon. The inner-membrane transporter MexB functions as a proton-motive-force-dependent antiporter, supported by the periplasmic membrane-fusion protein adaptor MexA and the outer-membrane duct OprM. Substrates captured from the inner leaflet or periplasm are extruded across the outer membrane through the trimeric OprM channel without releasing them into the cytoplasm.

\citet{cowan1999plant} catalogued plant secondary metabolites with antimicrobial activity. Phenolics and flavonoids were described as membrane-active. Tannins were described as protein-binding and metal-chelating. Terpenoids were described as lipophilic disruptors of the bilayer. Alkaloids, saponins, and quinones were assigned intercalation, pore formation, and thiol binding, respectively. The review is still the standard starting point for a crude-extract screen because it links solvent polarity to the metabolite class recovered.

\citet{hancock2000antibiotic} summarised why \taxa{Pseudomonas aeruginosa} fails many antimicrobial agents that inhibit other Gram-negative rods. Low outer-membrane permeability, inducible chromosomal AmpC $\beta$-lactamase, and multiple active RND efflux systems function concurrently to restrict intracellular drug concentrations. The paper is relevant here because one of the three test strains is \taxa{P. aeruginosa}, and non-linear zone responses across concentrations require evaluation against that envelope background.

\citet{masuda2000substrate} examined substrate specificities of the \taxa{P. aeruginosa} efflux systems MexAB-OprM, MexCD-OprJ, and MexXY-OprM. MexAB-OprM displayed broad substrate promiscuity, accommodating structurally diverse amphiphilic and lipophilic compounds, aromatic structures, and hydrophobic organic molecules.

\citet{nikaido2003molecular} restated the structural basis of Gram-negative outer-membrane permeability. Asymmetric lipopolysaccharide (LPS) packaging in the outer leaflet, stabilised by divalent cation bridges, retards passive permeation of hydrophobic terpenoids, whereas polar solutes depend upon proteinaceous water-filled porin channels. In crude aqueous extracts, polar phenolics and saponins encounter this outer leaflet as their initial cellular barrier.

\citet{poole2005efflux} reviewed active efflux as a coordinated Gram-negative resistance mechanism. RND pumps in \taxa{P. aeruginosa} operate synergistically with the low-permeability outer envelope: slow inward diffusion through narrow or gated porins enables active extrusion to maintain sub-inhibitory intracellular concentrations of diverse xenobiotics, including plant polyphenols and glycosides.

\citet{pages2008porin} described porin channel regulation as a selective diffusion barrier. Down-regulation or conformational closure of outer-membrane channels limits the influx of hydrophilic agents. In an agar diffusion assay, porin gating and active extrusion govern the threshold concentration required to arrest bacterial lawn growth.

\citet{zhao2011phytochemicals} comprehensively reviewed the phytochemical profile and biological activities of \textit{Dipsacus} species. Triterpenoid saponins and iridoids represent the primary specialized metabolite classes characteristic of the genus. Pentacyclic triterpenoid saponins of \textit{Dipsacus} are based on an oleanane skeleton, primarily featuring hederagenin ($3\beta,23$-dihydroxyolean-12-en-28-oic acid) as the aglycone. Among these, dipsacoside B is a prominent bidesmosidic saponin with a C-3 glycosidic chain (hederagenin 3-$O$-$\alpha$-L-rhamnopyranosyl-(1$\to$2)-$\alpha$-L-arabinopyranoside) and a C-28 ester-linked gentiobiose moiety (28-$O$-$\beta$-D-glucopyranosyl-(1$\to$6)-$\beta$-D-glucopyranosyl ester). Asperosaponin VI (akebia saponin D) carries an identical hederagenin core with an arabinopyranoside at C-3 and a gentiobiosyl ester at C-28, serving as the official statutory quality marker for \textit{Dipsaci Radix} in pharmacopoeial monographs. The amphiphilic character of these bidesmosidic saponins imparts surfactant properties, enabling them to perturb lipid packing and form transient pores in microbial bilayers. \citet{zhao2011phytochemicals} also documented monoterpene iridoids, notably loganin, sweroside, and cantleyoside derivatives, confirming that hot aqueous extractions recover both saponins and iridoids in abundance.

\citet{liebold2011growth} tested lipophilic extracts of \taxa{Dipsacus sylvestris} roots against \taxa{Borrelia burgdorferi} \textit{sensu stricto}. Growth was inhibited \invitro{}. The solvent was not water, and the organism was a spirochaete, not an enteric rod. The result still supports the broader claim that \textit{Dipsacus} tissues yield biologically active antimicrobial fractions.

\citet{azwanida2015review} compared maceration, Soxhlet, ultrasound, and related methods. Soxhlet was described as exhaustive and reproducible for a defined solvent, with heat as the main drawback for thermolabile constituents. Distilled water as the Soxhlet solvent is therefore a deliberate choice: it matches decoction chemistry and accepts some thermal cost.

\citet{balouiri2016methods} reviewed diffusion and dilution assays. Disc diffusion remains the Kirby--Bauer format for a standard antibiotic \citep{bauer1966antibiotic,clsi2024m100}. Zone diameter is semi-quantitative. It is not a minimum inhibitory concentration. The present study uses that diffusion readout only: zone of inhibition, measured inclusive of the disc.

The Angiosperm Phylogeny Group IV classification \citep{apgiv2016} updated familial circumscriptions for flowering plants, placing the order Dipsacales within the asterid clade Campanulids (euasterids II). Under APG IV, the historically segregated family Dipsacaceae Juss.\ is treated as subfamily Dipsacoideae Eaton within a broadly circumscribed Caprifoliaceae Juss.\ \textit{sensu lato}. This phylogenetic placement groups \textit{Dipsacus} alongside honeysuckles and scabious herbs, reflecting shared morphological traits and common biosynthetic pathways for iridoid glucosides and oleanane saponins.

\citet{chevalier2017porins} detailed the structure, function, and regulation of \taxa{P. aeruginosa} outer-membrane porins, explaining why its outer envelope restricts polar solute uptake far more than that of \taxa{E. coli}. The major porin, OprF, exists primarily (over 95\%) as a closed two-domain conformer with low single-channel conductance, whereas only a minor fraction (1--5\%) adopts an open 16-stranded $\beta$-barrel conformation. This structural gating accounts for the low baseline permeability of the \taxa{P. aeruginosa} outer membrane, setting a stringent barrier to passive influx of water-soluble phytochemicals.

\citet{oszmianski2020roots} reported phenolic profiles and biological activity of root and leaf extracts of \taxa{Dipsacus fullonum}. The paper adds a European teasel to the set of chemically studied species and shows that leaves, not only roots, carry phenolics.

\citet{jubair2021review} updated mechanisms of plant compounds against multidrug-resistant bacteria: membrane leakage, efflux inhibition, quorum-sensing interference, and biofilm suppression. Saponins disrupt outer and plasma membrane bilayers through surfactant action, whereas polar iridoids and phenolics can act concurrently to impair cellular homeostasis.

\citet{amrcollaborators2022global} estimated the 2019 global mortality associated with bacterial AMR. \taxa{E. coli}, \taxa{K. pneumoniae}, and \taxa{P. aeruginosa} were among the pathogens that contributed most. That ranking is the clinical reason for choosing the three strains.

\citet{skala2023dipsacus} reviewed specialised metabolites of \textit{Dipsacus} and \textit{Scabiosa}. The authors reaffirmed that oleanane-type triterpenoid saponins---predominantly dipsacoside B and asperosaponin VI---together with iridoid glucosides such as cantleyoside, sweroside, and loganin, represent the hallmark chemical profile of the genus. For \taxa{D. inermis}, they recorded the Kashmiri name wopal haakh, a Himalayan distribution, and traditional use against inflammation, cough, and related complaints. The review is the most compact recent statement of the genus-level chemistry.

The World Health Organization bacterial priority list of 2024 keeps carbapenem-resistant Enterobacterales and \taxa{P. aeruginosa} in the critical group \citep{who2024bacterial}. A student zone-of-inhibition screen does not address that clinical need directly. It can only ask whether a local aqueous extract still clears a lawn of those genera.

Soxhlet extraction with water is slower than a one-pot decoction, but it repeats the solvent until the thimble runs pale. That matters for tannins and saponins, which can take several siphons to leave the powder. Heat is the cost: some glycosides may hydrolyse. The present protocol accepts that cost because the research question is the aqueous extract, not a cold macerate \citep{azwanida2015review,harbone1998phytochemical}.

\citet{drakhshaan2026identification} is the most detailed antibacterial study of \taxa{D. inermis} to date. Dichloromethane and methanol extracts were tested by agar well diffusion at \SI{50}{\micro\gram\per\milli\litre} and \SI{100}{\micro\gram\per\milli\litre} against \taxa{Staphylococcus aureus}, \taxa{Bacillus subtilis}, \taxa{E. coli}, \taxa{P. aeruginosa}, \taxa{Salmonella typhi}, and \taxa{Enterobacter aerogenes}, with gentamicin \SI{10}{\micro\gram} discs as the reference. Dichloromethane at the higher concentration gave the largest zones, including \SI{15.93 \pm 0.11}{\milli\metre} against \taxa{E. coli} and \SI{16.83 \pm 0.29}{\milli\metre} against \taxa{P. aeruginosa}. Methanol was weaker. GC-MS of the dichloromethane fraction and HPLC of the methanol fraction pointed to terpenoids, steroids, and phenolics. \taxa{K. pneumoniae} was not tested, and water was not the extraction solvent. Those two omissions define the present study. Direct millimetre comparison with the present C1 and C2 values is not valid on a milligram-per-millilitre scale.

\section{National studies (India)}

\citet{ahmad2001antimicrobial} screened ethanolic extracts of forty-five Indian medicinal plants against drug-resistant clinical bacteria and \textit{Candida}. Forty extracts were active against one or more bacteria. Phenols, tannins, and flavonoids were the metabolites most often detected in the active fractions. The work established a national template: traditional use, solvent extract, diffusion assay, and a short phytochemical panel. The present dissertation follows that template with water rather than ethanol and with a Himalayan teasel rather than the north-Indian taxa they used.

\citet{laxminarayan2016antibiotic} described the drivers of antibiotic resistance in India: high infectious burden, over-the-counter access, hospital overuse, and limited diagnostics. Gram-negative rods in Indian hospitals were already difficult to treat. The paper supplies the national public-health frame for testing \taxa{E. coli}, \taxa{P. aeruginosa}, and \taxa{K. pneumoniae} in an Indian university laboratory.

\citet{taneja2019antimicrobial} reviewed AMR in the Indian environment as well as in the clinic. Resistant Enterobacterales and \taxa{Pseudomonas} circulate in water, waste, and farms as well as in wards. The paper also warns against reading an \invitro{} plant-extract zone as an environmental or therapeutic solution.

Indian laboratories routinely use Soxhlet extraction for pharmacognostic dissertations, with water, ethanol, or methanol as the solvent \citep{azwanida2015review}. Harborne's phytochemical handbook remains the source of the tube tests used here: Mayer's reagent for alkaloids, alkaline reagent for flavonoids, ferric chloride for phenols and tannins, froth for saponins, Salkowski for terpenoids and phytosterols, Fehling's for carbohydrates, and concentrated acid for quinones \citep{harbone1998phytochemical}. Those tests are qualitative. A colour or precipitate is not a concentration.

Indian hospital laboratories follow CLSI disc tables for gentamicin and related drugs \citep{clsi2024m100}. Those tables apply to the disc content, not to a plant extract. The GEN 5 disc in this study is a concurrent reference. It is not used to call the extract susceptible or resistant.

Indian teaching laboratories also publish Soxhlet yields and crude zone diameters without converting those diameters into broth MIC values \citep{azwanida2015review,balouiri2016methods}. That restraint is kept here. The antibacterial procedure is limited to disc diffusion for gentamicin GEN 5 and to zone-of-inhibition measurement, inclusive of the disc \citep{bauer1966antibiotic,balouiri2016methods}.

The present work stays inside that national method set: defined powder-to-solvent ratio, continuous hot extraction, concentration, and a short qualitative tube panel before the diffusion assay \citep{harbone1998phytochemical,ahmad2001antimicrobial}. What it adds is a Himalayan teasel and an aqueous solvent on three Gram-negative isolates held in a Kashmir zoology laboratory.

Taken together, the national literature supports three practical choices: aqueous Soxhlet extraction, a classical tube panel, and a diffusion assay scored as zone of inhibition, with GEN 5 as the disc reference. It also supports caution. Indian Gram-negative isolates are often harder to inhibit than laboratory strains used in older European papers.

\section{Regional studies (Jammu and Kashmir)}

\citet{lone2015study} documented medicinal plant use in Bandipora district, Kashmir. \taxa{Dipsacus inermis}, recorded as wopal hakh (wopal haakh), appears in that inventory of local remedies. The paper is an ethnobotanical catalogue, not an antibacterial assay. It is the earliest dated regional source used here that names the species in a Kashmiri district flora of use.

\citet{hassan2020dipsacus}, working at the University of Kashmir, showed that \taxa{D. inermis} modulates inflammation by inhibiting the NF-$\kappa$B pathway \invitro{} and \invivo{}. The study links a local plant to a defined signalling path. It does not report zones of inhibition. It does show that Himalayan material of this species is chemically active in a university laboratory in Srinagar, which is the same institutional setting as the present work. The anti-inflammatory end-points cannot be read as antibacterial end-points, and the present LB plates cannot be read as an NF-$\kappa$B assay.

\citet{bhat2021ethnobotany} reviewed medicinal plant use in Kashmir regions and argued for documentation before that knowledge is lost. Himalayan teasels and related herbs sit in that wider list of taxa used for inflammation, cough, and gastrointestinal complaints. The review does not isolate an aqueous antibacterial protocol.

\citet{haq2023cross} recorded high-altitude botanical resources among ethnic groups in the Kashmir Himalaya. Leaves of \taxa{D. inermis} were used for internal injuries and body inflammation. That is a specific folkloric claim, not a laboratory result. It is consistent with the anti-inflammatory work of \citet{hassan2020dipsacus} and with the saponin--phenolic chemistry summarised by \citet{skala2023dipsacus}.

\citet{ahad2023ethnobotanical} worked with Dard informants in Gurez. Crushed, boiled leaves of \taxa{D. inermis} (wopal haakh) were used as a postnatal bath, and the leaves were also cooked as food. Boiling in water is an aqueous preparation. The present Soxhlet water extract is a laboratory analogue of that solvent choice, not a reconstruction of the bath.

No regional paper located for this review reports an aqueous Soxhlet extract of Kashmir-grown \taxa{D. inermis} scored by zone of inhibition against \taxa{E. coli}, \taxa{P. aeruginosa}, and \taxa{K. pneumoniae} together. \citet{drakhshaan2026identification} remains the nearest antibacterial benchmark, but it used organic solvents and did not include \taxa{K. pneumoniae}. The present chapters fill that narrow regional and methodological gap without treating the result as a clinical recommendation.

\section{Synthesis}

From 1979 to 2026 the international literature progressed from iridoid and saponin structural characterisation \citep{jensen1979cantleyoside,cowan1999plant} through Gram-negative envelope mechanics and efflux architecture \citep{li1995role,hancock2000antibiotic,masuda2000substrate,nikaido2003molecular,chevalier2017porins} to genus-level phytochemical profiling of \textit{Dipsacus} \citep{zhao2011phytochemicals,skala2023dipsacus} under APG IV phylogenetic circumscription \citep{apgiv2016} and an organic-solvent antibacterial study of \taxa{D. inermis} \citep{drakhshaan2026identification}. Indian work supplied the AMR setting and the extraction-and-screen template \citep{ahmad2001antimicrobial,laxminarayan2016antibiotic,taneja2019antimicrobial}. Kashmir work supplied names, uses, and an anti-inflammatory mechanism for the same species \citep{lone2015study,hassan2020dipsacus,haq2023cross,ahad2023ethnobotanical}. The open experimental question is the aqueous Soxhlet extract, read as millimetre zones of inhibition on LB plates, with GEN 5 as the disc reference.

\FloatBarrier
